% TOP_PROBLEM: {"id": 2794, "title": "Kirby Problem 2.46", "category_id": 7, "category": {"id": 7, "name": "topology", "display_name": "Topology", "description": "Properties preserved under continuous deformations.", "slug": "topology", "order_index": 7, "created_at": "2026-07-31T15:26:25.671Z"}, "status": "open"}
% FIRST_POSTED: null
% ENTRY_KIND: statement_only
% Imported from: batch11.tex; UnsolvedMath ID 2794
\documentclass[11pt]{article}
\usepackage[margin=25mm]{geometry}
\usepackage{fontspec}
\setmainfont{Latin Modern Roman}
\setsansfont{Latin Modern Sans}
\usepackage{amsmath,amssymb,mathtools}
\usepackage{unicode-math}
\setmathfont{Latin Modern Math}
\usepackage{xurl,hyperref}
\setcounter{secnumdepth}{0}
\setlength{\parindent}{0pt}
\setlength{\parskip}{5pt}
\setlength{\emergencystretch}{3em}
\newcommand{\sref}[2]{\hyperlink{source:#1:#2}{[S#2]}}
\newcommand{\cataloguescope}[1]{\paragraph{Recorded target.}#1}
\begin{document}
% UnsolvedMath ID 2794; source catalogue code KP-2.46
\setcounter{section}{2793}
\section{Kirby Problem 2.46}\label{2794}
{\small\sffamily UnsolvedMath ID 2794 · Topology}
\subsection{Definitions and mathematical statement}

\paragraph{Shared notation.}$\mathbb N=\{1,2,\ldots\}$, and $\mathbb Z,\mathbb Q,\mathbb R,\mathbb C$ denote the integers, rationals, reals and complex numbers. Any other conventions are specified locally.


An oriented surface $S$ is of infinite type if it is not a compact finite-genus surface with finitely many punctures and boundary circles. Its mapping class group $\operatorname{Mod}(S)$ consists of orientation-preserving homeomorphisms modulo isotopy, fixing boundary pointwise. An end is an element of the inverse limit of the sets of components of complements of compact subsets. An end is planar if it has a genus-zero neighborhood; otherwise it is accumulated by genus. Give $G=\operatorname{Mod}(S)$ the quotient of the compact-open topology. A subset $A\subset G$ is coarsely bounded if it has finite diameter for every continuous left-invariant pseudometric on $G$, equivalently its orbits are bounded in all continuous isometric actions. A coarsely bounded generating set generates $G$ as an abstract group. For the canonical large-scale geometry, choose a symmetric coarsely bounded generating set $A$ containing an open neighborhood of the identity. Such word metrics represent the quasi-isometry class of compatible maximal left-invariant metrics. Then the word metric is
\[d_A(g,h)=\min\{m\ge0:g^{-1}h=a_1\cdots a_m,\ a_i\in A\}.\]
Hyperbolicity means that this word-metric group, or its Cayley graph, has uniformly thin geodesic triangles. Its quasi-isometry type is independent of the chosen coarsely bounded generating neighborhood. A quasi-isometry distorts distances by a multiplicative and additive bounded amount and has coarsely dense image. For the second question, consider orientable surfaces with finitely many ends, every end accumulated by genus; retain their compact boundary components, if present, in the boundary-fixing mapping class convention.
\paragraph{Statement recorded in the supplied source.}
Among coarsely boundedly generated mapping class groups of infinite-type surfaces, characterize those that are Gromov-hyperbolic in this word metric, including the bounded cases. For surfaces $S,S'$ in the stated class with respectively $n,m\ge2$ ends and $n\ne m$, determine whether $\operatorname{Mod}(S)$ and $\operatorname{Mod}(S')$ are quasi-isometric. \sref{2794}{2} \sref{2794}{3}

\subsection{Short English statement}
Classify hyperbolic mapping class groups under coarse bounded generation, and compare the large-scale geometry of groups for surfaces with different finite numbers of genus ends.
\subsection{Sources}
\begin{description}
\item[\hypertarget{source:2794:1}{S1}] ULAM.AI and the UnsolvedMath Contributors, UnsolvedMath: A Curated Collection of Open Mathematics Problems, record 2794 (KP-2.46). CC BY 4.0. \url{https://huggingface.co/datasets/ulamai/UnsolvedMath}.
\item[\hypertarget{source:2794:2}{S2}] R. İnanç Baykur, Robion C. Kirby and Daniel Ruberman, eds., \emph{K3: A New Problem List in Low-Dimensional Topology}, Mathematical Surveys and Monographs 295, American Mathematical Society (2026), Problem 2.46 and its remarks; Chapters 1 and 2 for the stated conventions. \url{https://math.berkeley.edu/sites/default/files/surv-295-ruberman-watermarked-author-pdf.pdf}.
\item[\hypertarget{source:2794:3}{S3}] Christian Rosendal, \emph{Large scale geometry of metrisable groups}, arXiv:1403.3106v1 (2014), Theorem 3 and Section 2.2 (maximal metrics and open generating sets with relative property (OB)). \url{https://arxiv.org/pdf/1403.3106v1}.
\end{description}
\label{end:2794}
\end{document}
